\section{Health System Implications and Responses}
\label{sec:health_implications}

\subsection{Direct Health Effects}
According to the WHO African Region Framework 2024-2033:

\begin{itemize}
    \item \textbf{Heat-Related Illness:}
    \begin{itemize}
        \item 47\% increase in heat-related emergency department visits
        \item Disproportionate impact on outdoor workers and urban populations
        \item Rising incidence of heat stroke and cardiovascular complications
    \end{itemize}
    
    \item \textbf{Respiratory Conditions:}
    \begin{itemize}
        \item 63\% increase in asthma exacerbations linked to air quality
        \item Rising prevalence of chronic respiratory conditions
        \item Increased burden on emergency medical services
    \end{itemize}
\end{itemize}

\subsection{Indirect Health Effects}
The Lancet Countdown 2024 identifies key indirect impacts:

\begin{itemize}
    \item \textbf{Mental Health:}
    \begin{itemize}
        \item 42\% increase in climate-related anxiety and depression
        \item Post-traumatic stress following extreme weather events
        \item Disruption of community support systems
    \end{itemize}
    
    \item \textbf{Nutritional Health:}
    \begin{itemize}
        \item 28\% increase in childhood stunting in affected regions
        \item Micronutrient deficiencies due to reduced food diversity
        \item Compromised maternal nutrition affecting birth outcomes
    \end{itemize}
\end{itemize}

\subsection{Health System Capacity}
WHO assessment of African health systems shows:

\begin{itemize}
    \item Only 28\% of facilities have climate-resilient infrastructure
    \item 52\% lack reliable access to essential services during extreme events
    \item Critical shortages in trained healthcare workers
\end{itemize}

\subsection{Disease Surveillance and Response}
Current capabilities include:

\begin{itemize}
    \item Limited early warning systems for climate-sensitive diseases
    \item Gaps in real-time disease surveillance
    \item Insufficient laboratory capacity for emerging pathogens
\end{itemize}

\begin{keymessage}
\textbf{Critical Health System Needs:}
\begin{itemize}
    \item Strengthen emergency response capabilities
    \item Enhance disease surveillance systems
    \item Improve climate-resilient infrastructure
    \item Build healthcare workforce capacity
\end{itemize}
\end{keymessage}

\subsection{Vulnerable Populations}
Health impact assessment reveals:

\begin{itemize}
    \item \textbf{Children:}
    \begin{itemize}
        \item 3.7x higher risk of climate-sensitive diseases
        \item Increased vulnerability to heat stress
        \item Long-term developmental impacts from malnutrition
    \end{itemize}
    
    \item \textbf{Elderly:}
    \begin{itemize}
        \item Limited physiological adaptation capacity
        \item Higher rates of chronic disease complications
        \item Reduced access to healthcare during extreme events
    \end{itemize}
    
    \item \textbf{Pregnant Women:}
    \begin{itemize}
        \item Increased risk of adverse pregnancy outcomes
        \item Higher susceptibility to vector-borne diseases
        \item Limited access to maternal healthcare during disasters
    \end{itemize}
\end{itemize}

\subsection{Health System Adaptation Strategies}
Current initiatives include:

\begin{itemize}
    \item Development of climate-resilient healthcare facilities
    \item Integration of climate considerations into health planning
    \item Capacity building for healthcare workers
    \item Enhanced emergency response protocols
\end{itemize}

\begin{recommendation}
\textbf{Priority Health System Interventions:}
\begin{enumerate}
    \item Implement comprehensive heat-health action plans
    \item Strengthen disease surveillance and early warning systems
    \item Develop climate-resilient health infrastructure
    \item Enhance capacity for emergency medical response
    \item Improve access to essential health services
\end{enumerate}
\end{recommendation}
